\scanpage{100}
\deleted{\unsure{provient|provenant} de $H^1(U,G)$, pourvu que $U$ soit \inserted{pris} assez petit.}

On utilisera un argument de spécialisation.

Or en \deleted{\unsure{vertu de}} \inserted{suivant} Cartier, des deux suites exactes
\[
\begin{aligned}
0\to\mathcal O_X\xrightarrow{F}\mathcal O_X\to\mathcal O_X/\mathcal O_X^p\to0\\
1\to\mathcal O_X^*\xrightarrow{F}\mathcal O_X^*\to\mathcal O_X^*/\mathcal O_X^{*p}\to1
\end{aligned}
\]
\margin{($X$ quelconque !,}

\deleted{\unsure{on ti|on tire}} on tire
\[
\begin{aligned}
H^1(X,\alpha_p)&=H^1(X,\mathcal O_X)^{1-F}\simeq H^0(X,\mathcal O_X/\mathcal O_X^p)\xrightarrow{df}H^0(X,\widetilde{\underline{\Omega}}^1_X)\\
\unsure{{}_pH^1|pH^1|H^1}(X,\mu_p)&=H^1(X,\mathcal O_X)^{1-F}\simeq H^0(X,\mathcal O_X/\mathcal O_X^*)\xrightarrow{\frac{df}{f}}H^0(X,\widetilde{\underline{\Omega}}^1_X)
\end{aligned}
\]
et si $X$ est \uline{normal} les deux flèches écrites sont injectives [\deleted{\unsure{leurs|tous}} \inserted{en tous cas,} les noyaux s'identifient, soit aux \uline{diviseurs additifs} sur $X$ annulés par $\unsure{\sqrt{p}|F}$, soit aux diviseurs multiplicatifs annulés par $F$].

Donc on est ramené au

\uline{Lemme} \ $H^0(X,\widetilde\Omega^1_{X/k})\longrightarrow H^0(Y_{\overline K},\widetilde\Omega^1_{Y_{\overline K}/\overline K})$ est injectif

La vérification est presque triviale \ldots

\uline{Corollaire} \ $\underline{\mathrm{Pic}}^\tau_{X_{\overline K}/\overline K}\longrightarrow\underline{\mathrm{Pic}}^\tau_{\deleted{\unsure{X}}Y_{\overline K}/\overline K}$ a un noyau qui est un $p$-\uline{groupe fini unipotent},
% slanted
\margin{purement infinitésimal}
[le fait qu'il soit unipotent résulte déjà de $\operatorname{prof.}X\geq2$, grâce au \deleted{Bertini} Lefschetz-Grauert-Bertini]. ``sans partie fixe''. Si $X$ est normal (pas néc. de $\dim\geq2$), \inserted{$\underline{\mathrm{Pic}}_{X/k}$ est propre sur $k$.} Le ``sans partie fixe'' provient de la normalité; elle implique que l'on a une \uline{immersion fermée} $\underline{\mathrm{Pic}}^\tau_{X/k}\longrightarrow\prod_i\underline{\mathrm{Pic}}_{Y_i/k}$, les $Y_i$ des sections hyperplanes égalt normales. Donc pour prouver
% archive sheet number: 99

